% Short running form; the full title is too wide for the head's right box.
\runninghead{The King's Reckoning and the Forgiven Debtor}
\properlane{debt-forgiven-from-the-heart}{The King's Reckoning and the Forgiven Debtor: Mercy Received Without Measure, Given From the Heart}
\label{sec:debt-forgiven-from-the-heart}

Heard from its Gospel, the Mass is a lesson in what the king's mercy requires.
The parable ends with a sentence on every hearer: \latin{Sic et Pater meus
caelestis faciet vobis, si non remiseritis unusquisque fratri suo de cordibus
vestris}. St John Chrysostom and St Augustine read it as the condition of a
pardon already given. Every hearer is the servant who owed ten thousand talents
and was forgiven beyond what he asked, and the only thing the king asks in
return is that the forgiven forgive. Chrysostom finds that the cruel servant
revoked his own pardon by seizing his fellow; Augustine hears God say that he
will call back what he remitted, as the covenant of the Lord's Prayer states;
both insist that the forgiveness be given from the heart and not with the lips
alone. Read with them, the rest of the formulary surrounds the parable with
prayers for the mercy it describes. The Introit's verse blesses the undefiled,
whom Augustine still calls debtors; the Collect asks God's continual mercy over
his family and its devotion in good works; the Epistle forbids the forgiven man
to make his brother the enemy; the Secret offers the sacrifice by which God
willed to be appeased; and the Postcommunion asks that what the mouth has
received the pure mind may follow, as the Gospel asks that what the lips forgive
the heart may forgive.

\subsection{The sentence that ends the Gospel}

The difficulty lies in the parable's last two verses. The king had said, \latin{omne
debitum dimisi tibi}, ``I forgave thee all the debt''; now, angry, he hands the
servant to the torturers \latin{quoadusque redderet universum debitum}, ``until
he paid all the debt''. The same whole debt that was forgiven is exacted again.
And Jesus applies the sentence to his hearers with the words ``So also shall my
heavenly Father do to you''. Two questions follow. What is the condition on which
the king's mercy stands, and how is it to be met? And can a pardon once given
return? If it can, forgiveness from God seems never secure; if it cannot, the
king's second sentence seems to contradict his first. The Fathers answer the
first question together and the second in words that differ.

\subsection{Ten thousand talents and a hundred pence}

Everything in the parable depends on the disproportion of the two debts: ten
thousand talents against a hundred pence. The hundred pence are \latin{centum
denarios}, and a denarius is the day's wage for which the householder of another
parable hires his labourers, ``a penny a day'' (Mt 20:2): the fellow servant owes
about a hundred days' labour, a real debt but one a working man could in time
repay. Augustine, who supposes the ten thousand talents to be ``at least ten
thousand sins'', says that a hundred denarii ``perhaps are a thousand asses'',
and asks, ``But what was this to ten thousand talents?'' (\work{Sermon} 83.3).
Chrysostom draws the lesson at once: ``Seest thou how great the difference
between sins against man and against God? As great as between ten thousand
talents, and a hundred pence, or rather even much more.'' The difference arises
``both from the difference of the persons, and the constant succession of our
sins. For when a man looks at us, we stand off and shrink from sinning: but when
God sees us every day, we do not forbear'' (\work{Homilies on Matthew} 61.1).
Augustine makes every hearer both debtors at once: ``every man is at once God's
debtor, and hath also some brother a debtor to himself. For who is there who is
not God's debtor, but he in whom there can be found no sin? And who is there who
hath not a brother his debtor, but he against whom no one hath sinned?''
(\work{Sermon} 83.2). St Jerome puts the argument in a single question, from the
greater to the lesser: if the lord so readily forgave his servant ten thousand
talents, \latin{quanto magis servi conservis suis debent minora dimittere?}
(\work{In Matthaeum} III, PL~26, col.~133).

Augustine goes further and finds the Law itself in the numbers. The Law is given
in ten precepts; the ten thousand talents and the hundred pence are both built
on ten, so that ``in both numbers you will find every kind of sin included''
(\work{Sermon} 83.6). He also reckons from St Luke's genealogy seventy-seven
generations from Christ back to Adam, the number of the Lord's answer to Peter
(Mt 18:22), which the Mass does not read. These
reckonings are Augustine's own. What they serve in his sermon is the point the
parable makes without them: there is no sin that
the forgiven may refuse to forgive. ``And so let us be ready to forgive all the
trespasses which are committed against us, if we desire to be forgiven'' (83.4).

\subsection{The threat that was mercy, and the gift beyond the request}

The king's first command, that the servant be sold with his wife and children,
looks cruel. Chrysostom says it was ``Not of cruelty, nor of inhumanity \dots\ but
of unspeakable tenderness. For it is His purpose to alarm him by this threat,
that He might bring him to supplication, not that he should be sold.'' The king
did not forgive before the reckoning because he wished the servant to learn
``from how many obligations He is delivering him, that in this way at least he
might become more mild towards his fellow servant'' (61.3). The servant asked
for little: ``Have patience with me, and I will pay thee all.'' He wanted time. And
the lord, moved with pity, gave him remission. ``Seest thou again surpassing
benevolence? The servant asked only for delay and putting off the time, but He
gave more than he asked, remission and forgiveness of the entire debt'' (61.3).
The servant's promise to pay all was empty; the lord did not hold him to it.

Then comes the scene the parable has prepared. The fellow servant falls down and
says the same words, \latin{Patientiam habe in me, et omnia reddam tibi}, and the
Gospel repeats them exactly. The forgiven servant hears his own prayer and
refuses it. Chrysostom marks the repetition: ``he did not regard even the words by
which he had been saved (for he himself on saying this was delivered from the ten
thousand talents), and did not recognize so much as the harbor by which he
escaped shipwreck'' (61.4). The fellow servants who saw it were grieved; ``They
therefore who did not owe, partook of the grief'' (61.4).

\subsection{Who revokes the pardon?}

The king's second sentence exacts the debt he had remitted. Who revoked the
pardon? Chrysostom puts the question to the servant himself as he seizes his
fellow: ``What doest thou, O man? perceivest thou not, thou art making the demand
upon thyself, \dots\ and revoking the sentence and the gift?'' When the lord then
pleads with the servant and excuses himself, ``being on the point of revoking his
gift'', Chrysostom corrects the phrase: ``or rather, it was not he that revoked
it, but the one who had received it''. The delivery to the torturers is ``for
ever; for he will never repay''. And Chrysostom notes what this means for the gifts of God,
which Paul calls ``without repentance'' (Rom 11:29): ``wickedness has had such
power as to set aside even this law. What then can be a more grievous thing than
to be revengeful, when it appears to overthrow such and so great a gift of God''
(61.4).

Augustine puts the same sentence in God's mouth: ``For I have forgiven thee
first; do thou at least forgive after that. For if thou wilt not forgive, I will
call thee back, and put upon thee again all that I had remitted to thee. For the
Truth doth not speak falsely; Christ neither deceiveth, nor is deceived''
(\work{Sermon} 83.7). For Chrysostom the receiver revokes the gift; for Augustine
God calls it back. The two formulations do not contradict each other: the one
looks at the cause, the servant's cruelty, and the other at the judge who
ratifies it. Jerome calls the closing sentence fearful for this reason:
\latin{Formidolosa sententia, si juxta nostram mentem sententia Dei flectitur
atque mutatur}, a fearful sentence, if God's sentence bends and changes according
to our own mind; and he says that it overturns every pretence of a feigned peace,
\latin{omnem simulationem fictae pacis evertit} (\work{In Matthaeum} III,
col.~133). St Rabanus Maurus, commenting on the parable, repeats Jerome's sentence
and adds one he gives under St Gregory's name: what was forgiven us through
penance is exacted of us again if we do not forgive from the heart what is done
against us (\work{Commentaria in Matthaeum}, on Mt 18:35).

St Thomas Aquinas takes up the question these formulations raise and answers it
more narrowly. The last verse, he says, seems to suggest that forgiven sins
return, as Origen held; but that cannot be so, since the remission of sins has its
efficacy from the sacraments. Forgiven sins are said to return \latin{per
ingratitudinem}, through ingratitude (\work{Super Matthaeum} 18). Thomas does not
deny what Chrysostom and Augustine say of the servant's fate; he denies that the
sins once remitted come back as the same sins: they are said to return through
ingratitude. The ungrateful man makes himself a debtor again, by a new sin, and
the parable sets that debt at the size of the gift refused. The hearer is not told that a grudge undoes his baptism; he is
told that to refuse mercy after receiving it is to make himself again a debtor
of the whole.

\subsection{From the heart}

The sentence ends with three words that the Fathers do not let pass: \latin{de
cordibus vestris}. Chrysostom says that the Lord requires two things here, ``both
to condemn ourselves for our sins, and to forgive others; and the former for the
sake of the latter, that this may become more easy (for he who considers his own
sins is more indulgent to his fellow-servant); and not merely to forgive with the
lips, but from the heart'' (61.5). Augustine: ``Say not with the tongue, `I
forgive,' and put off to forgive in the heart; for by His threat of vengeance God
showeth thee thy punishment. God knoweth where thou speakest. Man can hear thy
voice; God looketh into thy conscience.'' And he adds a sentence that reverses the
ordinary preference: ``Better is it that thou shouldest be violent in words, and
forgive in the heart, than in words be soft, and in the heart relentless''
(\work{Sermon} 83.7).

Augustine also shows what ``from the heart'' does not forbid. A boy who sins and
asks forgiveness, and sins again, and asks again, may still be corrected: ``Let us
reprove with words, and if need be with scourges; but let us withal forgive the
sin, and cast away the remembrance of it from the heart. For therefore did the
Lord add, `from your hearts,' that though through affection discipline be
exercised, gentleness might not depart out of the heart.'' The patient cries out
under the surgeon's knife, and is cut; ``This is not cruelty'', Augustine says of
the surgeon (83.8). Forgiveness from the heart
forbids rancour, not correction.

The same sermon names the place where every Christian makes this promise daily:
``Forgive us our debts, as we forgive our debtors.'' ``What debts of thine? All,
or a certain part? Thou wilt answer, All. So then do thou with thy debtor. This
then is the rule thou layest down, this the condition thou speakest of; this the
covenant and agreement thou dost mention when thou prayest'' (83.4). Bl.\ Ildefonso
Schuster calls the parable ``the most authoritative commentary on those words of
the Lord's Prayer'', and adds, ``The measure of divine compassion is fixed by
ourselves, and fully corresponds to that which we may have used towards our
neighbours'' (\work{The Sacramentary} III, p.~180).

\subsection{Who is the debtor?}

The Fathers disagree on one point of identification, and the disagreement shapes
the reading. Some interpreters, St Jerome reports, took the debtor of ten thousand
talents to be the devil; Origen had suggested as much in a higher sense he did not
fix. Jerome rejects it outright: \latin{Scio quosdam istum qui debebat decem
millia talenta, diabolum interpretari \dots\ nec Ecclesiasticae interpretationis
est, nec a prudentibus viris recipiendae} (col.~133), a reading neither of the
Church nor to be accepted by prudent men. Chrysostom and Augustine make the
debtor every hearer; Augustine says that ``every man begins from Baptism; he goes
out free, the `ten thousand talents' are forgiven him; and when he goes out, he
will soon find some fellow-servant his debtor'' (83.7). St Thomas makes the king
God and the servants every man, and he finds the title in Job: ``Hast thou
considered my servant Job?'' (Job 1:8). Rabanus, following an allegorical line of
his own, reads the debtor as the people of the Jews bound by the ten
commandments and the fellow servant as the Gentiles; he too then turns to every
hearer with Jerome. The moral identification is the one the parable's last
sentence demands, \latin{unusquisque fratri suo}: the debtor is each of us, and
the fellow servant each brother.

\subsection{The reckoning and its end}

\latin{Voluit rationem ponere}: the king willed to settle accounts. St Thomas
says that this reckoning begins when God first sends some tribulation and ends at
the day of judgment (\work{Super Matthaeum} 18; the third reading gives his
Latin); the liturgical commentators who expound the parable at its own place in their books
read it of the judgment too, Durandus saying that on that day the Lord will settle
accounts with us (\work{Rationale} VI.139), Honorius that Christ the true king will
then come and require of each his deeds (\work{Gemma animae} IV.93). St Alphonsus
Liguori, preaching on the verse \latin{tradidit eum tortoribus}, calls it ``a
description of the sentence of the eternal death'': those who die in sin, he
says, can make no satisfaction in the other life and remain debtors to the divine
justice for ever (\work{Sermons for All the Sundays}, Sermon L). He agrees with Chrysostom's ``for
ever; for he will never repay''. The continuation of \work{The Liturgical Year},
commenting on this Mass, sets beside the Gospel the petition of the \latin{Dies
irae}, ``grant us the gift of forgiveness, before the Day of reckoning cometh'',
and says: ``We are all of us, in fact, that negligent servant, that insolvent
debtor'' (XI, pp.~464--466).

\subsection{The whole Mass around the parable}

The other elements of the formulary pray for, warn of, or embody the mercy the
parable describes. The Introit confesses a king whose will holds all things;
the people for whom Mardochai prays is under sentence, and in lines the Mass does
not sing he asks, ``have mercy on thy people'' (Esth 13:15). The verse blesses
the undefiled who walk in the law, and Augustine, expounding it, will not let them
forget the apostle's ``If we say that we have no sin, we deceive ourselves''
(\work{Enarrationes} 118, sermon~1.3); Cassiodorus says that the saints become
undefiled \latin{quando indulgentiae munera consequuntur}, when they receive the
gifts of pardon. The undefiled of the Introit are forgiven debtors.

The Collect asks for the whole household the \latin{continua pietas} that the king
showed his servant, and asks that the family be devoted \latin{in bonis actibus}.
Augustine names the two works of mercy that answer the Gospel, ``Forgive, and ye
shall be forgiven: give, and it shall be given unto you'' (Lk 6:37--38; \work{Sermon}
83.2), and the good works the Collect asks include both. The Epistle, read beside
the parable, forbids the forgiven man to make his brother the enemy, and
Chrysostom himself joins the two: the man who has been wronged is to pour his
anger on the devil, not on the one who wronged him (\work{Homilies on Ephesians}
22). The Gradual names God the refuge of generation after generation, and
Augustine says he became our refuge ``ever since we turned to Thee''; the Alleluia
sings Israel released from Egypt, and Augustine adds that ``no man is turned away
in heart from this world, unless aided by the gift of the Divine mercy''. The
release of the Exodus and the release of the debtor are both the work of mercy.

The Offertory's Job is the man who bore injury without sin, and Chrysostom, in
the homily on this parable, sets him before the hearer as proof that even the
devil's malice can become a cause of crowns, so that no man need be feared as an
enemy (61.5; quoted in the first reading). St Thomas, expounding the servant's ``Have patience with
me'', sets beside it a cry of Job's (Job 10:20). The Secret offers the sacrifice
by which God willed to be appeased, \latin{quibus et te placari voluisti}; in the
parable the lord was appeased by a prayer, and Chrysostom reminds the servant of
``God, whom by mere prayer thou hast reconciled'' (61.4). Schuster reads the
Secret's \latin{potenti pietate} as the mercy that ``has triumphed over divine
justice itself'' (p.~181).

The Communion's plea for judgment on persecutors, sung after this Gospel, would
seem to undo it. Augustine answers in two places. Of anger he says: ``A man
lusteth to be avenged, and Christ is not yet avenged, the holy martyrs are not yet
avenged. Still doth the patience of God wait, that the enemies of Christ, the
enemies of the martyrs, may be converted. \dots\ Forgive therefore; from the heart
forgive'' (\work{Sermon} 58.8). And of the martyrs who in the Apocalypse cry to
be avenged he asks ``whether they pleaded against the men themselves or against
the dominion of sin'', and answers that ``it is a genuine avenging of the martyrs,
and one full of righteousness and mercy, that the dominion of sin should be
overthrown''. That dominion is overthrown partly by men's amendment and partly by
the condemnation of the incurable; and he points to St Paul, who in bringing his
own body into subjection avenged the martyr Stephen ``in his own person''
(\work{De sermone Domini in monte} I.22.77). The plea for judgment and the
command to forgive ask the same thing, the end of sin's rule, and the first way
it ends is the persecutor's conversion.

The Postcommunion asks that what the mouth has received, \latin{quod ore
percepimus}, the pure mind may follow. Chrysostom's ``not merely \dots\ with the
lips, but from the heart'' and Augustine's ``Say not with the tongue \dots\ and put
off to forgive in the heart'' are said of the Gospel's forgiveness. The prayer
asks the same agreement of mouth and mind of the one who has received the
sacrament, and so the Mass ends where the parable ends, at the heart.

\subsection{The strongest difficulty}

The reading has to hold two things together. The first is the parable's
severity. Augustine and Chrysostom do not soften it: the debt returns, God calls it
back, the delivery to the torturers is for ever. St Thomas's distinction does not
weaken that severity but places it: sins once remitted do not return as the same
sins, but ingratitude makes a man a debtor again. Thomas tells the hearer that a
grudge does not undo his baptism; Augustine and Chrysostom tell him that refusing
mercy after receiving it puts the whole debt back on him. Both are true of the
parable, and each keeps the other from being misheard.

The second is the reach of ``from the heart''. The command could be taken to
forbid every correction, every refusal, every justice. Augustine's sermon refuses
that reading: discipline may be exercised ``through affection'', as the surgeon
cuts to cure, so long as gentleness does not depart from the heart. What the
command forbids is the remembrance kept against the brother, the forgiveness of
the lips that hides a refusal of the heart, what Jerome calls \latin{simulatio
fictae pacis}.

\subsection{The four senses of this reading}

\begin{fourSenses}
\item[Literal.] Jesus tells his disciples a parable: a king forgives a servant an
unpayable debt on his plea for time; the servant refuses a fellow servant a small
one on the same plea; the king calls him wicked and hands him to the torturers
until he pays all; so will the heavenly Father do to those who do not forgive their
brothers from the heart. In Matthew the parable answers Peter's question how often
to forgive (18:21--22), which this Mass does not read.
\item[Allegorical.] The king is God (Thomas), and in Rabanus's allegory Christ; the
ten thousand talents are the sins forgiven every man in baptism (Augustine); the
fellow servants who grieve and tell their lord are the members who suffer with
one member that suffers (Thomas); the
remission given on the servant's prayer is the Church's daily ``Forgive us our
debts, as we forgive our debtors'', which Augustine calls a covenant.
\item[Moral.] Forgive every debt, as one forgiven far more, and forgive it from the
heart, not with the lips alone (Chrysostom, Augustine); correct when correction is
owed, but without rancour (Augustine); remember your own sins, and the brother's
will seem small.
\item[Anagogical.] The king's reckoning, begun in tribulation, is completed at the
judgment (Thomas; Durandus; Honorius), when God's measure will be the measure we
used (Schuster); the debtor who would not forgive remains a debtor for ever
(Chrysostom; Alphonsus); and the Church asks before that day, in the words Fromage
sets beside this Gospel, for ``the gift of forgiveness, before the Day of
reckoning''.
\end{fourSenses}
